\section{Framework: tools, context, and memory}
\label{sec:framework}
We define research capability as the ability to choose appropriate quantitative methods,
interpret their outputs, and adapt future choices based on realized outcomes.
Counting API calls alone does not measure this capability. Our framework separates
tool execution, within-task context selection, and cross-episode memory.

\paragraph{Decision and learning cycle.}
Let $q_t$ be the task prompt, $\mathcal{D}_{\le t}$ the market data available at time $t$,
and $M_t$ the persistent memory state. An attention routing mechanism prunes interaction
history to fit the context budget. The agent uses the pruned context and a readout of
$M_t$ to select a tool, execute it with chosen parameters, and produce an asset
allocation. When the trade clears at settlement $t+T$, the realized return updates
$M_t$, provided it passes validation checks.

\begin{table*}[t]
\centering\small
\setlength{\tabcolsep}{5pt}
\begin{tabular}{p{0.16\textwidth}p{0.26\textwidth}p{0.25\textwidth}p{0.23\textwidth}}
\toprule
Component & Role in research workflow & Target question & Implementation and validation \\
\midrule
\textsc{fin-skills} & Document conventions and execute quantitative algorithms &
Which methods does the model select, and how are outputs applied? &
129 skills, 48 backends; evaluated on 64 portfolio decisions \\
HiSTrim & Prune historical KV cache while preserving task instructions &
Which interaction tokens can be dropped without harming decisions? &
Dynamic routing with sign preservation; evaluated on sequential tasks \\
Persistent memory & Update associative weights from completed episode returns &
Does past experience improve future method selection across regimes? &
Multi-timescale mushroom-body model; evaluated on delayed return series \\
Execution checks & Bind verification to code artifacts and filter invalid feedback &
Are execution outputs and learning signals free of look-ahead leaks? &
36 guards; evaluated against 12 planted defect families \\
\bottomrule
\end{tabular}
\caption{System components, roles, and validation scope within the financial agent framework.}
\label{tab:framework}
\end{table*}

\paragraph{Two distinct memory mechanisms.}
HiSTrim operates within a single session by evicting tokens from the KV cache during
generation. Persistent memory operates across sessions by storing parameter updates in an
external module. The two mechanisms operate on different timescales: context pruning
controls what the model attends to right now, whereas persistent memory retains summary
statistics of what happened in previous episodes.

\paragraph{Grounded feedback updates.}
When an episode concludes, the agent records the market context, the chosen tool, and the
realized return. Only verified, settlement-cleared outcomes update memory. Negative returns
are retained as learning signals, while trades relying on look-ahead data or unadjusted
prices are rejected by guards to prevent memory corruption.
